\begin{notechapter}
\chapter{Fully Convolutional Networks and Upsampling}\label{ch:fully-convolutional}
\chaptersubtitle{Networks for inputs of any size, dense predictions, and transposed convolutions}

\begin{chapterintro}
A convolutional layer does not care about the size of its input: the same kernels slide over
a small or a large image. Only fully connected layers fix the input size. A network made of
convolutions alone, a \emph{fully convolutional network}, accepts inputs of any size and can
produce outputs of matching size: a prediction per pixel, or per cell of a grid. This
chapter shows how to build such networks, how to enlarge the receptive field without
losing resolution, and how to upsample with transposed convolutions. You will learn
\begin{itemize}[nosep]
  \item how global pooling and $1\times1$ convolutions make a classifier size-agnostic, in code;
  \item dilated convolutions, whose receptive field grows exponentially with depth;
  \item a complete example: the actor--critic network of a lawnmower coverage agent, which
        runs on lawns of any size;
  \item the transposed convolution as the transpose of a convolution's matrix, as a scatter
        operation, and as zero insertion; its output size and its checkerboard artefacts;
  \item encoder--decoder networks for semantic segmentation: FCN and U-Net.
\end{itemize}
\prerequisites{\cref{ch:convolutional-layers,ch:cnn-architectures}.}
\end{chapterintro}

\section{Removing the fixed input size}\label{fully-convolutional:sec:anysize}

A fully connected layer has a weight matrix of fixed shape, so a network ending with
\texttt{Flatten()} and \texttt{Linear(64*8*8, 128)} (\cref{convolutional-layers:sec:network}) only
accepts $32\times32$ images. Two changes remove the restriction:
\begin{itemize}
  \item replace the fully connected classifier by a $1\times1$ convolution, which computes class scores
        at every position of the last feature map;
  \item reduce the score map to one vector by \emph{global average pooling} (or global max pooling),
        whose output has one value per channel whatever the spatial size (Lin, Chen and Yan,
        \emph{Network in Network}, ICLR, 2014).
\end{itemize}
By \cref{convolutional-layers:exr:5}, the $1\times1$ convolution followed by averaging equals averaging
followed by a linear layer: the classifier is a linear model on the average feature vector.

\begin{lstlisting}[language=Python]
import torch
from torch import nn

class TinyFCN(nn.Module):
    """Classifier that accepts images of any size >= 4x4."""
    def __init__(self, n_classes=10):
        super().__init__()
        self.features = nn.Sequential(
            nn.Conv2d(3, 32, 3, padding=1), nn.ReLU(), nn.MaxPool2d(2),
            nn.Conv2d(32, 64, 3, padding=1), nn.ReLU(), nn.MaxPool2d(2),
            nn.Conv2d(64, n_classes, 1))          # 1x1 conv: class scores per position
    def forward(self, x):
        scores = self.features(x)                 # (B, n_classes, H/4, W/4)
        return scores.mean(dim=(2, 3))            # global average pooling -> (B, n_classes)

net = TinyFCN()
for h, w in [(32, 32), (57, 91), (224, 224)]:
    x = torch.randn(1, 3, h, w)
    print((h, w), tuple(net.features(x).shape), tuple(net(x).shape))
# (32, 32) (1, 10, 8, 8) (1, 10)
# (57, 91) (1, 10, 14, 22) (1, 10)
# (224, 224) (1, 10, 56, 56) (1, 10)
\end{lstlisting}

The same weights produce a $10$-vector for every input size; the intermediate score map has
the size of the input divided by $4$ (rounded down by the poolings). Each of its positions is
the classifier applied to one receptive field ($10\times10$ pixels,
\cref{convolutional-layers:sec:receptive}).

\paragraph{Convolutionalization.} A trained classifier with fully connected layers can be
converted. VGG's first fully connected layer acts on a $512\times7\times7$ map; it is exactly a $7\times7$
convolution with $4096$ output channels and no padding, which produces $1\times1$ on a $224\times224$ input
and a larger map on a larger input. The next fully connected layers become $1\times1$ convolutions.
The converted network evaluates the classifier at all positions of a sliding window in
one pass, sharing the computation of overlapping windows (Long, Shelhamer and Darrell,
\emph{Fully Convolutional Networks for Semantic Segmentation}, CVPR, 2015).

\section{Large receptive fields without downsampling: dilation}\label{fully-convolutional:sec:dilation}

Poolings and strides enlarge the receptive field but discard resolution, which matters
when the output must be precise at the pixel level. A \emph{dilated} (or atrous) convolution with
dilation $d$ applies its $k\times k$ weights to pixels $d$ apart (\cref{convolutional-layers:def:conv}): a $3\times3$ kernel
with $d = 2$ covers $5\times5$ pixels, with $9$ weights. By the recursion of
\cref{convolutional-layers:sec:receptive} with stride $1$, each such layer adds $2d$ to the receptive
field. With dilations $1, 2, 4,\dots,2^{L-1}$,
\[
  r_L = 1 + 2\sum_{l=0}^{L-1}2^l = 2^{L+1} - 1:
\]
the receptive field grows exponentially with depth, at full resolution and with a linear
number of parameters (Yu and Koltun, \emph{Multi-Scale Context Aggregation by Dilated
Convolutions}, ICLR, 2016). With padding $p = d$ the spatial size is unchanged. A stack of
equal dilations only samples a sparse grid of the input (``gridding''); increasing dilations,
or mixing in undilated layers, avoids it.

\section{Case study: a coverage agent for lawns of any size}\label{fully-convolutional:sec:lawnmower}

The \texttt{Lawnmower} project (\url{https://github.com/CL-BZH/lawnmower}) trains an
actor--critic agent (trained by PPO) to steer a mower over a grid lawn until every cell is
mowed. Lawns differ in size and shape, and a policy trained on small lawns should run on
large ones. The network is therefore fully convolutional. Its observation has $4$ channels
on an $H\times W$ grid (cells the mower cannot reach, mowed cells, the mower's position, and the cell
ahead of it), and it outputs $8$ action logits (the moves to the neighbouring cells) and a value.

\begin{lstlisting}[language=Python]
class ResBlock(nn.Module):
    def __init__(self, ch, dilation):
        super().__init__()
        self.conv1 = nn.Conv2d(ch, ch, 3, padding=dilation, dilation=dilation)
        self.conv2 = nn.Conv2d(ch, ch, 3, padding=1)

    def forward(self, x):
        return x + self.conv2(F.relu(self.conv1(F.relu(x))))

class CoverageNet(nn.Module):        # abridged from lawnmower/learn/net.py
    def __init__(self, c_in=4, c=32, dilations=(1, 2, 4, 8, 16), hidden=128):
        super().__init__()
        self.stem = nn.Conv2d(c_in, c, 3, padding=1)
        self.trunk = nn.Sequential(*(ResBlock(c, d) for d in dilations))
        self.policy = nn.Sequential(nn.Linear(9 * c, hidden), nn.ReLU(),
                                    nn.Linear(hidden, 8))
        self.value = nn.Sequential(nn.Linear(3 * c + 1, hidden), nn.ReLU(),
                                   nn.Linear(hidden, 1))

    def forward(self, obs, mask=None):   # obs: (B, 4, H, W), any H, W
        feat = F.relu(self.trunk(self.stem(obs)))   # (B, c, H, W)
        # policy: features of the 3x3 cells around the mower -> 8 logits
        # value: mower features, mean and max over the lawn, cells left
        ...
\end{lstlisting}

\begin{figure}[ht]
\centering
\begin{tikzpicture}[x=1cm, y=1cm]
  \def\nnbaseline{-1.2}
  \nnbox{in}{0}{0.15}{2.0}{1.6}{gray}{observation\\$4\times H\times W$}
  \nnbox{st}{1.9}{0.35}{2.0}{1.6}{bookaccent}{stem\\$32\times H\times W$}
  \nnbox{r1}{3.5}{0.35}{2.0}{1.6}{bookaccent}{res.\\$d=1$}
  \nnbox{r2}{4.9}{0.35}{2.0}{1.6}{bookaccent}{res.\\$d=2$}
  \nnbox{r3}{6.3}{0.35}{2.0}{1.6}{bookaccent}{res.\\$d=4$}
  \nnbox{r4}{7.7}{0.35}{2.0}{1.6}{bookaccent}{res.\\$d=8$}
  \nnbox{r5}{9.1}{0.35}{2.0}{1.6}{bookaccent}{res.\\$d=16$}
  \nnarrow{in}{st}\nnarrow{st}{r1}\nnarrow{r1}{r2}\nnarrow{r2}{r3}\nnarrow{r3}{r4}\nnarrow{r4}{r5}
  % policy and value heads
  \node[draw, rounded corners=2pt, fill=bookexample!20, font=\scriptsize, align=center] (pol) at (12.6,1.0)
    {$3\times3$ patch at mower\\$9\cdot32\to128\to8$ logits};
  \node[draw, rounded corners=2pt, fill=bookdef!20, font=\scriptsize, align=center] (val) at (12.6,-1.0)
    {mower, mean, max, left\\$97\to128\to1$ value};
  \draw[-{Stealth}, thick, gray!70!black] (r5-out) -- (pol.west);
  \draw[-{Stealth}, thick, gray!70!black] (r5-out) -- (val.west);
\end{tikzpicture}
\caption{The lawnmower agent's network. The trunk keeps the full resolution (no pooling,
no stride: the exact position of the mower and of each unmowed cell matters); dilations
$1,2,\dots,16$ give each cell a receptive field of $75\times75$ cells. The heads read a fixed number of
features, gathered at the mower's position or pooled over the lawn, so the whole network
works for every lawn size.}
\label{fully-convolutional:fig:lawnmower}
\end{figure}

Three ideas make it size-agnostic:
\begin{itemize}
  \item \textbf{A resolution-preserving trunk.} The stem and the five residual blocks use
        padding equal to the dilation, so the feature map stays $H\times W$, one $32$-vector per cell.
        The receptive field after the stem is $3$, and each block adds $2d$ (dilated layer) $+ 2$ (undilated
        layer): $r = 3 + 2(1 + 2 + 4 + 8 + 16) + 5\cdot2 = 75$ cells, which the gradient of a central feature with
        respect to the input confirms. Each cell's features summarize a $75\times75$ neighbourhood.
  \item \textbf{A policy read at the mower.} The policy head takes the $9$ feature vectors of the mower's
        cell and its eight neighbours (zero-padded at the border) and outputs one logit per
        move; illegal moves are masked by setting their logits to $-10^9$. Its input size, $9\times32$, does not
        depend on the lawn.
  \item \textbf{A value from global pooling.} The value head takes the mower's features, the mean
        and the maximum of the features over the lawn cells (masked pooling, ignoring
        obstacles), and the number of cells left: $3\times32 + 1 = 97$ numbers, whatever the lawn.
\end{itemize}
The same $144\,361$ parameters run on a $10\times12$ and on a $40\times40$ lawn. On lawns wider than $75$ cells, a
cell's features no longer see the whole lawn, but the pooled features of the value head
still do.

\section{Upsampling}\label{fully-convolutional:sec:upsampling}

Dense prediction networks first reduce the resolution to gain context, then must
return to the input resolution. Fixed upsampling (nearest neighbour, bilinear
interpolation) has no parameters; the \emph{transposed convolution} is a learned upsampling.

\paragraph{The transpose of a convolution.} A strided convolution is a linear map
$y = Cx$ (\cref{convolution-filtering:sec:definitions}). In one dimension, with kernel $w = (w_0, w_1, w_2)$, stride $2$
and no padding, from $\R^5$ to $\R^2$:
\[
  C = \begin{bmatrix} w_0 & w_1 & w_2 & 0 & 0\\ 0 & 0 & w_0 & w_1 & w_2\end{bmatrix},\qquad
  C^\top\begin{bmatrix}y_0\\y_1\end{bmatrix} = \begin{bmatrix}w_0y_0\\ w_1y_0\\ w_2y_0 + w_0y_1\\ w_1y_1\\ w_2y_1\end{bmatrix} .
\]

\begin{definition}[Transposed convolution]\label{fully-convolutional:def:transposed}
The transposed convolution with kernel $w$, stride $s$ and padding $p$ is the linear map $C^\top$, where $C$ is
the matrix of the convolution with the same kernel, stride and padding. It maps the output
space of $C$ back to its input space.
\end{definition}

So it is the transpose of the convolution, \emph{values included}, not its inverse: $C^\top C\neq I$ in
general. It is also exactly the backward pass of the convolution with respect to its input
(\cref{convolutional-layers:prop:backward}): frameworks implement both with the same
routine. With $w = (1,2,3)$ and $y = (1,-1)$, $C^\top y = (1, 2, 2, -2, -3)$, which \texttt{conv\_transpose1d} returns.

\paragraph{Scatter view.} Reading $C^\top y$ column by column: each input value $y_i$ is multiplied by the
whole kernel, and the result is added into the output at position $si$; overlapping
contributions add up. A convolution \emph{gathers} a window into one value; its transpose
\emph{scatters} one value into a window. Padding $p$ in the corresponding convolution means that $p$
values are cropped from each end of the scattered result. Hence the output size
\[
  n_{\mathrm{out}} = (n - 1)s - 2p + k\ (+\ \text{output\_padding}),
\]
the inverse of the convolution's size formula; since the convolution rounds down, several
input sizes map to the same output size, and \texttt{output\_padding} selects one.

\begin{example}[A $2\times2$ map upsampled to $3\times3$]\label{fully-convolutional:ex:transposed}
Input $\begin{bmatrix}a&b\\c&d\end{bmatrix}$, kernel $W = (w_{uv})_{u,v=1}^3$, stride $2$, padding $1$
(\texttt{ConvTranspose2d(1, 1, 3, stride=2, padding=1)}). Scattering the four scaled kernels at
offsets $(0,0)$, $(0,2)$, $(2,0)$, $(2,2)$ gives a $5\times5$ map
\[
  \begin{bmatrix}
    aw_{11} & aw_{12} & aw_{13} + bw_{11} & bw_{12} & bw_{13}\\
    aw_{21} & aw_{22} & aw_{23} + bw_{21} & bw_{22} & bw_{23}\\
    aw_{31} + cw_{11} & aw_{32} + cw_{12} & aw_{33} + bw_{31} + cw_{13} + dw_{11} & bw_{32} + dw_{12} & bw_{33} + dw_{13}\\
    cw_{21} & cw_{22} & cw_{23} + dw_{21} & dw_{22} & dw_{23}\\
    cw_{31} & cw_{32} & cw_{33} + dw_{31} & dw_{32} & dw_{33}
  \end{bmatrix},
\]
and cropping one row and column on each side gives the $3\times3$ output
\[
  \begin{bmatrix}
    w_{22}a & w_{23}a + w_{21}b & w_{22}b\\
    w_{32}a + w_{12}c & w_{33}a + w_{31}b + w_{13}c + w_{11}d & w_{32}b + w_{12}d\\
    w_{22}c & w_{23}c + w_{21}d & w_{22}d
  \end{bmatrix}.
\]
With $w_{uv} = 3(u - 1) + v$ (the numbers $1$ to $9$ row by row) and $(a,b,c,d) = (2,3,5,7)$ this is
$\begin{bmatrix}10&24&15\\26&61&38\\25&58&35\end{bmatrix}$, as PyTorch computes.
\end{example}

\paragraph{Zero insertion.} Equivalently: insert $s - 1$ zeros between neighbouring input values,
pad with $k - 1 - p$ zeros on each side, and apply an ordinary stride-$1$ cross-correlation with the
kernel rotated by $180^\circ$. The inserted zeros explain the name \emph{fractionally strided convolution}
(Dumoulin and Visin, \emph{A Guide to Convolution Arithmetic for Deep Learning}, 2016).
``Deconvolution'', also used, is a misnomer: nothing is inverted.

\paragraph{Checkerboard artefacts.} When the stride does not divide the kernel size,
output positions receive different numbers of contributions. With $k = 3$, $s = 2$, a constant
input and a constant kernel, the interior output alternates between $1$, $2$ and $4$ contributions per
pixel, a checkerboard pattern that training has to cancel and often does not (Odena,
Dumoulin and Olah, \emph{Deconvolution and Checkerboard Artifacts}, Distill, 2016). Kernel
sizes divisible by the stride ($k = s = 2$, as below) avoid the uneven overlap; the alternative is
\emph{resize-convolution}: nearest or bilinear upsampling followed by an ordinary convolution.

\section{Semantic segmentation}\label{fully-convolutional:sec:segmentation}

Semantic segmentation assigns a class to every pixel. A fully convolutional network
outputs a map of class scores $Z\in\R^{K\times H\times W}$, and the loss is the cross-entropy summed over pixels,
$-\sum_{i,j}\log\softmax(Z[:,i,j])_{y_{ij}}$ (\cref{ch:sigmoid-softmax}).

\paragraph{FCN.} Long et al.\ (2015) convolutionalized a classifier, upsampled its coarse
score map (stride $32$) with transposed convolutions, and added the score maps predicted from
earlier, finer layers (strides $16$ and $8$): \emph{skip connections} that recover the detail lost by
pooling.

\paragraph{U-Net} (Ronneberger, Fischer and Brox, \emph{U-Net: Convolutional Networks for
Biomedical Image Segmentation}, MICCAI, 2015) makes the idea symmetric: a contracting
encoder, an expanding decoder with transposed convolutions, and at each resolution the
encoder's feature map \emph{concatenated} to the decoder's (\cref{fully-convolutional:fig:unet}). Encoder
features say \emph{where}, decoder features say \emph{what}. U-Nets are the standard for
segmentation with little data, and the backbone of diffusion models.

\begin{figure}[ht]
\centering
\begin{tikzpicture}[>=Stealth, x=1cm, y=1cm, lab/.style={font=\scriptsize}]
  % encoder
  \fill[bookaccent!35] (0,-1) rectangle (0.3,1);      \draw[bookaccent!70!black] (0,-1) rectangle (0.3,1);
  \fill[bookaccent!35] (1.6,-3.0) rectangle (2.0,-1.8); \draw[bookaccent!70!black] (1.6,-3.0) rectangle (2.0,-1.8);
  \fill[bookaccent!35] (3.6,-4.55) rectangle (4.2,-3.85); \draw[bookaccent!70!black] (3.6,-4.55) rectangle (4.2,-3.85);
  % decoder: upsampled (green) + concatenated skip, then convs (blue)
  \fill[bookdef!35] (5.8,-3.0) rectangle (6.2,-1.8);   \draw[bookdef!70!black] (5.8,-3.0) rectangle (6.2,-1.8);
  \fill[bookaccent!35] (6.3,-3.0) rectangle (6.7,-1.8); \draw[bookaccent!70!black] (6.3,-3.0) rectangle (6.7,-1.8);
  \fill[bookdef!35] (7.8,-1) rectangle (8.1,1);        \draw[bookdef!70!black] (7.8,-1) rectangle (8.1,1);
  \fill[bookaccent!35] (8.2,-1) rectangle (8.5,1);     \draw[bookaccent!70!black] (8.2,-1) rectangle (8.5,1);
  \fill[bookexample!40] (9.6,-1) rectangle (9.9,1);    \draw[bookexample!70!black] (9.6,-1) rectangle (9.9,1);
  % arrows
  \draw[->, thick, bookpitfall] (0.15,-1) -- (1.8,-1.8) node[midway, above right, lab] {pool};
  \draw[->, thick, bookpitfall] (1.8,-3.0) -- (3.9,-3.85) node[midway, above right, lab] {pool};
  \draw[->, thick, bookdef] (3.9,-3.85) -- (6.0,-3.0) node[midway, right, lab] {up};
  \draw[->, thick, bookdef] (6.5,-1.8) -- (7.95,-1) node[midway, right, lab] {up};
  \draw[->, thick, bookexample, dashed] (0.3,0.5) -- (7.8,0.5) node[midway, above, lab] {copy and concatenate};
  \draw[->, thick, bookexample, dashed] (2.0,-2.1) -- (5.8,-2.1) node[midway, above, lab] {copy and concatenate};
  \draw[->, thick] (8.5,0) -- (9.6,0) node[midway, above, lab] {$1\times1$};
  % sizes
  \node[lab, anchor=east] at (-0.05,0) {$c\times H\times W$};
  \node[lab, anchor=east] at (1.55,-2.6) {$2c\times\frac H2\times\frac W2$};
  \node[lab, anchor=east] at (3.55,-4.2) {$4c\times\frac H4\times\frac W4$};
  \node[lab, anchor=west] at (9.95,0) {$K\times H\times W$};
\end{tikzpicture}
\caption{A two-level U-Net. Blue: two $3\times3$ convolutions with ReLU; green: upsampled
features (transposed convolution, $k = s = 2$) concatenated with the copied encoder features
(dashed); orange: per-pixel class scores.}
\label{fully-convolutional:fig:unet}
\end{figure}

\begin{lstlisting}[language=Python]
def block(c_in, c_out):
    return nn.Sequential(nn.Conv2d(c_in, c_out, 3, padding=1), nn.ReLU(),
                         nn.Conv2d(c_out, c_out, 3, padding=1), nn.ReLU())

class TinyUNet(nn.Module):
    def __init__(self, c_in=3, n_classes=2, c=16):
        super().__init__()
        self.enc1, self.enc2 = block(c_in, c), block(c, 2 * c)
        self.mid = block(2 * c, 4 * c)
        self.pool = nn.MaxPool2d(2)
        self.up2 = nn.ConvTranspose2d(4 * c, 2 * c, 2, stride=2)
        self.dec2 = block(4 * c, 2 * c)            # 2c (skip) + 2c (up)
        self.up1 = nn.ConvTranspose2d(2 * c, c, 2, stride=2)
        self.dec1 = block(2 * c, c)
        self.head = nn.Conv2d(c, n_classes, 1)     # per-pixel class scores

    def forward(self, x):                          # H, W multiples of 4
        e1 = self.enc1(x)                          # c   x H   x W
        e2 = self.enc2(self.pool(e1))              # 2c  x H/2 x W/2
        m = self.mid(self.pool(e2))                # 4c  x H/4 x W/4
        d2 = self.dec2(torch.cat([e2, self.up2(m)], dim=1))
        d1 = self.dec1(torch.cat([e1, self.up1(d2)], dim=1))
        return self.head(d1)                       # n_classes x H x W
\end{lstlisting}

The network returns a $2\times64\times64$ map for a $64\times64$ image and a $2\times96\times160$ map for a $96\times160$ image. The
constraint ``multiple of $4$'' comes from the two poolings: the upsampled maps must match the
encoder maps they are concatenated with.

\begin{chaptersummary}
\begin{itemize}[nosep]
  \item Only fully connected layers fix the input size. A $1\times1$ convolution and global pooling make a
        classifier size-agnostic; fully connected layers can be rewritten as convolutions.
  \item Dilations $1,2,\dots,2^{L-1}$ give a receptive field $2^{L+1} - 1$ at full resolution.
  \item A fully convolutional trunk plus heads that gather features at chosen positions or pool
        them globally works on grids of any size, as in the lawnmower agent.
  \item The transposed convolution is $C^\top$, the transpose of the convolution's matrix: a scatter of
        scaled kernels, or zero insertion plus a convolution with the rotated kernel; it is the
        convolution's input gradient. Uneven overlaps cause checkerboards.
  \item FCN and U-Net combine coarse semantic features with fine encoder features through skip
        connections to predict a class per pixel.
\end{itemize}
\end{chaptersummary}

\section*{Exercises}
\addcontentsline{toc}{section}{Exercises}

\begin{exercise}\label{fully-convolutional:exr:1}
Compute the shape of the score map of \texttt{TinyFCN} for a $33\times35$ input. What is the smallest input
it accepts?
\solutionlink{fully-convolutional:exr:1}
\end{exercise}

\begin{exercise}\label{fully-convolutional:exr:2}
Compute the receptive field of a stack of $3\times3$ convolutions with dilations $1, 2, 4, 8, 1, 2, 4, 8$, and the
number of layers of undilated $3\times3$ convolutions needed for the same receptive field.
\solutionlink{fully-convolutional:exr:2}
\end{exercise}

\begin{exercise}\label{fully-convolutional:exr:3}
For the matrix $C$ of \cref{fully-convolutional:sec:upsampling} with $w = (1,2,3)$, compute $C^\top C$ and check that it is not
the identity. Which input vectors $x$ satisfy $Cx = 0$?
\solutionlink{fully-convolutional:exr:3}
\end{exercise}

\begin{exercise}\label{fully-convolutional:exr:4}
Compute the output size of \texttt{ConvTranspose2d} with $k = 4$, $s = 2$, $p = 1$ on a $16\times16$ input. Does this
kernel produce a checkerboard on a constant input? Compare with $k = 3$, $s = 2$.
\solutionlink{fully-convolutional:exr:4}
\end{exercise}

\begin{exercise}\label{fully-convolutional:exr:5}
In the lawnmower network, what would the receptive field be with dilations $(1,1,1,1,1)$? With
$(1,2,4,8,16,32)$?
\solutionlink{fully-convolutional:exr:5}
\end{exercise}
\end{notechapter}
